\subsection{Experimental Setup}\label{sec:results-setup}

\paragraph{Training data.}
Models are trained on the CT part of
TotalSegmentator~\cite{wasserthalTotalSegmentatorRobustSegmentation2023}
(1082 training, 57 validation, and 89 test subjects) and, except in
\autoref{sec:results-fusion}, on the TotalSegmentator MRI release~\cite{dantonoliTotalSegmentatorMRIRobust2025}
(561 training subjects from University Hospital Basel and the NIH Imaging
Data Commons). Training uses 61 of the 117 classes, chosen to balance
anatomical categories and structure thickness: 12 abdominal and pelvic organs,
8 thoracic, head, and spinal organs, 5 muscles (one side of each pair), 6 limb,
shoulder, and pelvic bones, 11 vertebrae spanning all spine regions, 9 ribs
and sternal structures, and 10 vessels. The remaining 56 classes form the
unseen split. (\autoref{tab:totalseg-classes}) It includes structures close to training classes, such as the
contralateral side of each training muscle. On MRI, 31 of the 61 training
classes are available.

\paragraph{Model configurations.}
The experiments were run at different stages of development, so the three
results sections use different model configurations
(\autoref{tab:configs}). Numbers from different sections are therefore not
directly comparable. In particular, TotalSegmentator MRI is a held-out
modality for the fusion experiments, which train on CT only, but part of the
training data in the other two sections.

\begin{table}[t]
  \centering
  \small
  \caption{Model configurations used in the results sections.}
  \label{tab:configs}
  \begin{tabular}{@{}lllll@{}}
    \toprule
    Experiment & Section & Training data & Levels & $p_\text{synth}$  \\
    \midrule
    Image--Label Fusion & \ref{sec:results-fusion} & TotalSeg CT & 1 & 0  \\
    Cascade & \ref{sec:results-cascade}, \ref{sec:results-synth} & TotalSeg CT + MRI & 2--3 & 0 \\
    Synth & \ref{sec:results-synth} & TotalSeg CT + MRI & 1 & 0.5 \\
    Synth + Cascade & \ref{sec:results-synth} & TotalSeg CT + MRI & 2--3 & 0.5 \\
    \bottomrule
  \end{tabular}
\end{table}
\todo[inline]{Resolve before submission: report a single final configuration,
or keep this split?}

\paragraph{Baselines.}
The external baseline is Medverse~\cite{huMedverseUniversalModel2025a} with
its released weights (\autoref{sec:results-cascade},
\autoref{sec:results-synth}). In \autoref{sec:results-fusion}, we additionally
train Medverse's architecture from its released weights with our training
recipe.
\todo[inline]{No Iris baseline exists yet; the comparison in
\autoref{sec:fundamentals-relatedwork} remains architectural, not empirical.}

\paragraph{Evaluation protocol.}
Each target is segmented with a single context pair ($K = 1$). The context
subject is drawn at random, with a fixed seed per target, from the same dataset
and split as the target, among the subjects that contain the target class.
Every subject that contains a class serves once as a target for that class. For each model, we use the checkpoint with the best
Dice on the in-distribution validation split; no evaluation dataset is used for
checkpoint selection.
\todo[inline]{See open points: MRI validation split, $n$ in
\autoref{tab:ood-sources}, $K$ during training.}


\paragraph{Metrics.}
Accuracy is reported as Dice, computed per class and averaged over classes
unless stated otherwise, and as Normalized Surface Dice (NSD) with a tolerance
of 3\,mm. Compute is reported as GFLOPs per forward pass and as latency per
sample. All training and evaluation runs on a single NVIDIA
RTX PRO 6000 Blackwell Server Edition GPU (96\,GB).
\todo[inline]{Precision of the latency measurements in
\autoref{tab:cascade-medverse} (fp32 as in \autoref{tab:fusion-staircase}?).}

\paragraph{Datasets.}
\autoref{tab:ood-sources} lists the evaluation datasets, grouped by their
relation to the training classes. TotalSegmentator and FLARE22 contain classes
from the TotalSegmentator vocabulary, so their classes split into seen and
unseen ones. The other datasets contain targets that do not occur in training:
new anatomical structures, lesions, a measurement region that is not an
anatomical object, and, for the fusion experiments, an unseen modality.
Because of this split, "held-out" is used at two levels below: at the
dataset level (a dataset whose images were never trained on, e.g.\ FLARE22,
whose classes are nonetheless mostly in-vocabulary) and, where stated
explicitly as such, at the class level (excluding classes that are in the
training vocabulary, even if their dataset is held out).

\begin{table}[t]
  \centering
  \small
  \caption{Evaluation datasets, grouped by relation to the training classes.
  $n$: number of subjects; $n_\text{cls}$: number of evaluated classes.}
  \label{tab:ood-sources}
  \begin{tabular}{@{}lllll@{}}
    \toprule
    Dataset & Modality & Target & $n$ & $n_\text{cls}$ (seen/unseen) \\
    \midrule
    \multicolumn{5}{@{}l}{\itshape Training vocabulary (seen and unseen classes)} \\
    TotalSegmentator~\cite{wasserthalTotalSegmentatorRobustSegmentation2023} & CT &  full-body 
    anatomical structures & 89 & 117 (61/56) \\
    FLARE22~\cite{maFLARE22Challenge2024} & CT & abdominal organs & 50 & 13 (11/2) \\
    \midrule
    \multicolumn{5}{@{}l}{\itshape New targets} \\
    ISLES22~\cite{hernandezpetzscheIsles2022Multicenter2022} & MRI DWI & stroke lesion & 250 & 1 \\
    Shifts-MS~\cite{malininShifts20Extending2022} & MRI FLAIR & MS lesions & 46 & 1 \\
    MSD Hippocampus~\cite{antonelliMedicalSegmentationDecathlon2022} & MRI T1 & hippocampus & 260 & 2 \\
    MSD Prostate~\cite{antonelliMedicalSegmentationDecathlon2022} & MRI T2 + ADC & prostate zones & 64 & 2 \\
    ATLAS v2.0$^{*}$~\cite{liewLargeCuratedOpensource2022} & MRI T1 & chronic stroke lesion & 654 & 1 \\
    GNC\_705 & MRI Dixon & kidney lesions & 610 & 13 \\
    NasalSeg~\cite{zhangNasalSegDataset2024} & CT & nasal cavity and sinuses & 107 & 5 \\
    \midrule
    \multicolumn{5}{@{}l}{\itshape Non-anatomical target} \\
    HU\_LWK1 & CT & center of L1 vertebrae & 36 & 1 \\
    \midrule
    \multicolumn{5}{@{}l}{\itshape Additional modality (unseen in \autoref{sec:results-fusion})} \\
    TotalSeg MRI~\cite{dantonoliTotalSegmentatorMRIRobust2025} & MRI & full-body 
    anatomical structures & 47 & 50 (31/19) \\
    \bottomrule
  \end{tabular}
  \par\vspace{0.5em}
  \raggedright\footnotesize
  $^{*}$Obtained from an unofficial mirror; provenance not yet confirmed.
\end{table}
